\documentclass[11pt,fontset=windows]{ctexart}

\usepackage[a4paper,margin=22mm]{geometry}
\usepackage{booktabs}
\usepackage{siunitx}
\usepackage{pgfplots}
\usepackage{xcolor}
\usepackage{hyperref}
\usepackage{caption}

\pgfplotsset{compat=1.18}
\sisetup{
  round-mode=places,
  round-precision=3,
  scientific-notation=true,
  table-number-alignment=center
}

\title{SSSS Mindlin 方板自由震動無因次頻率報告}
\author{Timoshenko / Mindlin 2D SSSS Q4 5x5}
\date{\today}

\begin{document}
\maketitle

\section{目的}

本報告整理 \texttt{2d\_ssss\_vibration\_mindlin\_exact\_comparison.csv} 中的自由震動結果，
用無因次頻率
\[
  \Omega = \omega a^2\sqrt{\frac{\rho h}{D}},
  \qquad
  D = \frac{Eh^3}{12(1-\nu^2)}
\]
比較 FEM 數值解與 SSSS Mindlin/Reissner 厚板解析解。使用 \(\Omega\) 的原因是它消除了尺寸、密度與彎曲剛度的量綱影響，比直接比較 Hz 更適合作為 benchmark 指標。

\section{無因次頻率折線圖}

圖~\ref{fig:omega-lines} 比較每個 mode rank 的 FEM 無因次頻率 \(\Omega_{\mathrm{FEM}}\) 與解析無因次頻率 \(\Omega_{\mathrm{exact}}\)。縱軸採用對數尺度，因為第 60 以後的 FEM 模態出現非常大的跳升。

\begin{figure}[htbp]
\centering
\begin{tikzpicture}
\begin{semilogyaxis}[
  width=0.92\linewidth,
  height=0.42\linewidth,
  grid=both,
  xlabel={Mode rank},
  ylabel={\(\Omega\)},
  legend pos=north west,
  legend cell align=left,
  tick align=outside,
]
\addplot+[blue, mark=*, mark size=1.2pt, line width=0.8pt]
  table[col sep=comma, x=mode_rank, y=Omega_num]
  {../data/2d_ssss/2d_ssss_vibration_mindlin_exact_comparison.csv};
\addlegendentry{\(\Omega_{\mathrm{FEM}}\)}

\addplot+[black, dashed, mark=square*, mark size=1.0pt, line width=0.8pt]
  table[col sep=comma, x=mode_rank, y=Omega_exact]
  {../data/2d_ssss/2d_ssss_vibration_mindlin_exact_comparison.csv};
\addlegendentry{\(\Omega_{\mathrm{exact}}\)}
\end{semilogyaxis}
\end{tikzpicture}
\caption{FEM 與 SSSS Mindlin 解析解的無因次頻率對比。}
\label{fig:omega-lines}
\end{figure}

\section{無因次頻率誤差折線圖}

圖~\ref{fig:omega-error} 顯示相對誤差百分比：
\[
  e_\Omega =
  \frac{\Omega_{\mathrm{FEM}}-\Omega_{\mathrm{exact}}}{\Omega_{\mathrm{exact}}}
  \times 100\%.
\]
第 60 以後誤差非常大，因此同樣使用對數尺度表示誤差絕對值。

\begin{figure}[htbp]
\centering
\begin{tikzpicture}
\begin{semilogyaxis}[
  width=0.92\linewidth,
  height=0.42\linewidth,
  grid=both,
  xlabel={Mode rank},
  ylabel={\(|e_\Omega|\) (\%)},
  legend pos=north west,
  legend cell align=left,
  tick align=outside,
]
\addplot+[red!75!black, mark=*, mark size=1.2pt, line width=0.8pt]
  table[col sep=comma, x=mode_rank, y expr=abs(\thisrow{rel_error_percent})]
  {../data/2d_ssss/2d_ssss_vibration_mindlin_exact_comparison.csv};
\addlegendentry{\(|\Omega_{\mathrm{FEM}}-\Omega_{\mathrm{exact}}|/\Omega_{\mathrm{exact}}\)}
\end{semilogyaxis}
\end{tikzpicture}
\caption{無因次頻率相對誤差百分比。}
\label{fig:omega-error}
\end{figure}

\section{指定模態數據表}

表~\ref{tab:selected-modes} 摘錄 mode 1、5、10、20、30、40、50、60、70、75 的無因次頻率與誤差。最後一欄中 \(1\) 表示特徵方程相對殘差低於程式設定容許值，\(0\) 表示未通過。

\begin{table}[htbp]
\centering
\caption{指定 mode rank 的 \(\Omega\) 對比表。}
\label{tab:selected-modes}
\begin{tabular}{
  r
  S[table-format=7.3]
  S[table-format=3.3]
  S[table-format=7.3]
  r
}
\toprule
{Mode} & {\(\Omega_{\mathrm{FEM}}\)} & {\(\Omega_{\mathrm{exact}}\)} & {Error (\%)} & {Residual OK} \\
\midrule
1  & 33.150      & 19.065  & 73.877      & 0 \\
5  & 186.811     & 85.038  & 119.679     & 1 \\
10 & 648.074     & 133.621 & 385.008     & 1 \\
20 & 664.914     & 220.600 & 201.411     & 1 \\
30 & 706.744     & 283.911 & 148.931     & 1 \\
40 & 753.931     & 351.899 & 114.246     & 1 \\
50 & 808.681     & 397.675 & 103.352     & 1 \\
60 & 3869883.724 & 440.013 & 879393.775  & 1 \\
70 & 3909987.223 & 492.178 & 794325.089  & 1 \\
75 & 5529556.951 & 507.629 & 1089191.159 & 1 \\
\bottomrule
\end{tabular}
\end{table}

\section{簡要解讀}

前 50 個模態中，FEM 的 \(\Omega\) 整體高於 Mindlin 解析解，代表目前 5x5 Q4 離散結果偏硬。第 60 以後 FEM \(\Omega\) 產生數個量級的跳升，這些模態不宜直接作為主要物理彎曲模態解讀。殘差欄位只表示離散特徵值問題 \(K\phi=\omega^2M\phi\) 的代入精度，不能取代與解析 \(\Omega\) 的比較。

\end{document}
